% Section 3: Verification Regimes as Taxonomic Move
%
% Status: first-pass integrated draft. Lending primary, underwriting as turn
% near the end. Section 2 took the unusual example (BSA/AML); this section
% takes the canonical one (lending) for load-management reasons articulated in
% planning document. Will be reworked.
%
% Citations needed: Rudin 2019, Bordt et al. 2022, Lanham et al. 2023, Turpin
% et al. 2023, ECOA / Reg B adverse action notice requirements, FCRA notice
% provisions, recent CFPB guidance on adverse action notices for AI-driven
% decisions, fair lending literature on disparate impact in pricing, recent FTC
% and CFPB statements on algorithmic pricing.

\section{Verification Regimes as Taxonomic Move}
\label{sec:verification-regimes}

Section~\ref{sec:frame-as-method} demonstrated that the empty-chair frame separates a governance objective, the artifact offered as evidence, and the conclusion drawn from that artifact. This section applies that distinction to consumer lending. AI applications in regulated decision contexts differ structurally according to the claims their verification regimes can support: which evidentiary objects the architecture exposes, how those objects bind to the decision, and what an affected party or examiner can test.

The worked example is the adverse action notice regime under the Equal Credit Opportunity Act and Regulation~B. The notice communicates specific reasons for a denial to an applicant who was absent from the decision but is legally entitled to an account of it. That function creates a precise verification question: what does the notice establish about the basis of the creditor's action, and what evidence allows the applicant or an examiner to test that relationship?

\subsection{The structural line}
\label{sec:verification-regimes:structural-line}

Recall the opening's framing: a verification regime lets a party test a material claim about a decision. Seven evidentiary objects recur in AI-supported workflows. Decision records preserve \emph{inputs and context}. SHAP can characterize \emph{functional feature dependence} under its value-function and background construction. LIME offers a \emph{local behavioral approximation} under its perturbation and surrogate choices \parencite{lundberg2017,ribeiro2016}. Attention or activation analyses expose selected \emph{internal computational states} \parencite{jain2019}. Chain-of-thought and other natural-language accounts are \emph{generated rationales}. Contemporaneous workflow records preserve observable parts of a \emph{human deliberative record}. Signatures and approval logs establish \emph{institutional authorization and review}. A system may expose several of these objects, but none acquires the evidentiary function of another merely because both are called explanations.

The methods therefore require method-specific readings. SHAP attributes an output under a specified construction; it does not by itself recover a human deliberative record or prove a complete causal account of model-internal computation. LIME approximates behavior locally, with a warrant bounded by its neighborhood and surrogate. An attention map exposes selected internal quantities, but attention weight alone is not a complete account of computational causation. Generated chain-of-thought is text whose faithfulness must be tested rather than assumed \parencite{lanham2023,turpin2023}. An ante-hoc interpretable model can expose its decision function more directly, while leaving separate questions about data provenance, policy validity, and deployment. A contemporaneous human record can show what was presented, consulted, recorded, and decided, but not every cognitive process. An institutional attestation proves that an authorized actor made a recorded commitment; whether the attested content is true depends on its evidence.

This taxonomy yields a structural line, but not a division between ``explanations'' and ``no explanations.'' An artifact cannot, by itself, verify a proposition about a different evidentiary object. Additional detail of the same kind does not close that gap. The gap closes when the architecture supplies an inspectable relation among the decision, the relevant object, the method that produced the artifact, and the proposition being tested. Section~\ref{sec:primitives} develops governance requirements for preserving and inspecting those relations. It does not promise transparent access to an anthropomorphic ``model reasoning'' object.

\subsection{Adverse action notices as the canonical case}
\label{sec:verification-regimes:adverse-action-notices}

The Equal Credit Opportunity Act requires creditors to provide adverse action notices to applicants who are denied credit, stating the specific reasons for the denial \parencite{regB1002_9}. The Fair Credit Reporting Act extends this to credit-based denials more broadly \parencite{fcraAdverseAction}. The regulatory intent is straightforward: a customer denied credit is entitled to know why, both to enable challenge of the denial and to inform the customer's understanding of their own creditworthiness. The empty chair the regime addresses is the denied applicant.

For a rule-based or low-dimensional scorecard, an institution may be able to trace a denial directly to the operative rule, threshold, input, and policy provision. Complexity changes the evidence problem, but it does not make every adverse action notice derived from a complex model unfaithful. Institutions may use feature attribution, local approximation, counterfactual analysis, model-specific decision paths, or combinations of methods. The CFPB does not permit complexity to excuse inaccurate or non-specific reasons \parencite{cfpbCircular2022_03}; the creditor must determine which method supports the reasons it communicates.

The notice itself establishes that the creditor communicated stated reasons. A method record can establish how those reasons were generated. Validation evidence can establish how the method behaved under specified tests. The further proposition --- that the communicated reasons accurately identify the factors material to this creditor's action on this application --- depends on the relation among those artifacts. Understandability is relevant but not sufficient; method fidelity is relevant but not sufficient; a plausible list is not the same as a tested binding to the action.

The empty-chair reading makes the evidentiary sequence visible. The applicant receives a legally required artifact. Their interest is represented substantively when the creditor can show why the artifact supports the proposition the law requires and when the applicant or examiner has a tractable way to challenge that showing. The frame does not declare every post-hoc method inadequate for every purpose. It asks the institution to name the purpose, the artifact, its warrant, and its limits.

The produced-absence dimension, picking up the distinction from Section~\ref{sec:frame-as-method}: the applicant who would have challenged the denial if challenge were tractable does not challenge. Adverse action notices are written in regulatory boilerplate, the reasons are typically generic (``credit history,'' ``debt-to-income ratio,'' ``insufficient income for amount requested''), and the applicant has no mechanism to interrogate whether these reasons accurately describe the model's behavior on their application \parencite{cfpbCircular2022_03}. The architecture produces silence (applicants who do not appeal because appeal is structurally unavailable) and reads the silence as evidence that the notice was adequate. The institution's position that adverse action notices satisfy ECOA / Reg~B is sustained partly by the fact that customers cannot effectively contest the notices' faithfulness, not by the notices being faithful.

\subsection{The taxonomic move}
\label{sec:verification-regimes:taxonomic-move}

The taxonomy divides applications by the verification claim their architecture can support.

\emph{AI that surfaces material for an independently accountable disposition.} Pattern recognition can surface applications for investigation; anomaly detection can flag unusual facts; consistency monitoring can identify departures from institutional norms. If an authorized reviewer receives the information and authority needed to make the disposition, the institution can preserve the decision inputs and context, the reviewer's observable actions, the recorded basis for the disposition, and the authorization. This does not make the human record infallible or expose every cognitive process. It produces an accountable institutional decision whose evidentiary components can be inspected.

\emph{AI that supplies the disposition and an account of it.} An underwriting model may generate a recommendation or decision, an explanation method may generate principal reasons, and a human may approve the result. This architecture is not unverifiable merely because it uses AI or a post-hoc method. Its burden is different: the institution must validate the method appropriate to its target, bind the resulting artifact to the operative model and case, define what the human review establishes, and preserve evidence for the broader substantive claim. A signature cannot supply a missing method-to-decision binding; nor does the presence of an explanation make such a binding impossible.

These are architecturally different, not simply points on a risk scale. The first locates the disposition and its recorded basis in an authorized human process supported by AI. The second locates material parts of the disposition in a model-supported process and therefore requires evidence across model behavior, explanation method, case context, policy, and review. ``Authority to override'' answers who may intervene. It does not answer what the reviewer could test, whether independent judgment was exercised, or what the approval establishes.

\subsection{What the categoricals foreclose, and what they do not}
\label{sec:verification-regimes:foreclosures}

The structural line forecloses specific \emph{claims}, not entire model classes. The following claims are not established without additional binding evidence:

\begin{itemize}
\item That an adverse action notice is faithful to the creditor's action solely because a post-hoc method generated plausible principal reasons and a human approved the notice.

\item That a decision-recommendation workflow received substantive human verification solely because an authorized reviewer signed or had formal override power.

\item That a retrospective justification reconstructs the decision-time information and deliberative record solely because no contradictory contemporaneous record exists.
\end{itemize}

The structural line permits stronger claims where the corresponding evidence is present, including:

\begin{itemize}
\item Ante-hoc interpretable models whose operative decision function, inputs, thresholds, and policy relationship are available for inspection, with separate validation of provenance and deployment.

\item Decision-support tools where AI surfaces material, an authorized human makes the disposition, and the architecture preserves the information presented, observable review activity, recorded basis, and authorization.

\item Model-supported decision architectures that validate an explanation method for its stated target, bind its output to the operative model and case, document its assumptions and limits, and preserve the other evidence needed for the substantive claim under review.
\end{itemize}

\subsection{A turn: pricing as produced-absence canonical case}
\label{sec:verification-regimes:pricing-turn}

The lending example demonstrates the categoricals at their canonical regulatory location, where the dishonest middle takes the form of an artifact (the adverse action notice) that misrepresents the reasoning behind a denial. This is not the only form the dishonest middle takes. We turn briefly to risk-based pricing in consumer lending (the determination of the interest rate, fees, and terms offered to an \emph{approved} applicant) to demonstrate that the frame produces interpretive leverage in regulatory spaces where the empty chair is not addressed by an artifact at all.

The Truth in Lending Act requires disclosure of a rate's mathematical structure (APR, finance charges, payment schedule); it does not require disclosure of why this particular rate was offered to this particular applicant \parencite{regZ1026}. Unlike the denied applicant, the priced applicant has no regulatory entitlement to a substantive explanation of their rate. The applicant receives a rate; the applicant accepts or declines; the institution proceeds.

The dishonest middle in pricing is more diffuse than in denial because there is no regulatory artifact equivalent to the adverse action notice. The pricing decision is not communicated to the customer as a reasoned determination; it is communicated as a take-it-or-leave-it offer. The customer accepts because they need credit and have no basis for negotiation. The institution reads the customer's acceptance as evidence that the rate was acceptable~\parencite{bartlett2022consumer}.

This is produced absence at architectural scale. The customer who would have challenged the rate if challenge were a thing they expected to receive does not challenge, because the architecture of risk-based pricing does not contemplate challenge as a normal part of the transaction. The institution reads the customer's silence (accepted offers, completed transactions, regulatory examinations finding no fair-lending issues) as evidence that pricing is fair. The architecture has produced the silence; the institution interprets the silence as endorsement.

The frame's diagnostic move on this configuration is the produced-absence move from Section~\ref{sec:frame-as-method}: \emph{what is the architecture doing that creates this silence, and what is it inferring from the silence it created?}

The architecture is producing pricing decisions through models whose reasoning is not communicated to the customer; structuring the transaction such that the customer cannot effectively negotiate without comparable offers they cannot tractably obtain; and producing aggregate pricing data that institutions analyze for fair-lending compliance, while the data itself is censored: customers who would have asked for justification if justification were available do not exist as a population the institution measures.

The institution then infers from this silence that pricing is fair. The customers who accepted are taken as evidence the rates were acceptable. The customers who declined are taken as having self-selected out. The customers who never applied because they expected unfavorable rates are not measured at all. The architecture has produced a population of accepted-and-silent customers and reads their silence as evidence of pricing fairness.

A common rejoinder: pricing models are statistically validated against outcomes (default rates by pricing tier, fair-lending disparate-impact analysis, ongoing monitoring \parencite{sr11_7,sr26_2}). This is true and important, and it is not a verification regime for the empty chair we are discussing. Statistical validation against outcomes verifies that the model performs as intended in aggregate; it does not verify that any individual customer's rate reflects their creditworthiness rather than model artifact. The aggregate empty chair (the regulator examining fair-lending compliance) is partly represented; the individual empty chair (the priced applicant) is not. The frame's contribution here is the aggregate-versus-individual distinction the regulatory regime currently elides.

\subsection{The taxonomic move, summarized}
\label{sec:verification-regimes:summary}

The lending example demonstrates the categoricals at the regulatory location where the dishonest middle is most directly addressable: where an artifact exists, where the regulation requires the artifact to represent the empty chair, and where post-hoc explanation methods have produced an artifact that does not. The pricing turn demonstrates that the frame extends to regulatory locations where the empty chair has been produced-absent by the regulatory regime itself, and that the frame proposes architectural primitives that would represent the chair the regulation does not currently require representing.

The remaining sections develop the diagnostic techniques that operationalize this taxonomic move (Section~\ref{sec:diagnostics}) and the categories of architectural primitive that empty-chair representation requires (Section~\ref{sec:primitives}). The frame's claim is that an honest AI governance framework must hold the structural line, both where the regulation requires it (as in adverse action notices) and where the regulation has not yet recognized that the line exists (as in risk-based pricing). Whether this claim is sustainable in practice is the subject of Section~\ref{sec:uncertainties}'s honest uncertainties.
